Throughout this study, we consider that our methods were effective in supporting the analysis of refactoring practices related to Type Hint adoption and software modernization in Python projects. However, several limitations may influence the interpretation of the results.

In this work, we operationalize the concept of software rejuvenation as refactoring activities related to the adoption of Type Hints. Although this definition captures an important aspect of modernization in the Python ecosystem, it represents only a partial view of the broader concept of software rejuvenation, not covering other forms of code evolution and improvement.

In addition, the dataset is restricted to open-source Python projects hosted on GitHub, selected based on activity and maturity criteria. While this enables the analysis of relevant and actively maintained systems, it may introduce a bias toward more established repositories, limiting generalization to smaller projects, other programming languages, or industrial contexts. The data selection process may also introduce sampling bias, since we focused on commits with a large number of detected transformations and selected one representative commit per developer, which may overrepresent large-scale modernization efforts and underrepresent incremental changes.

Regarding data collection, we relied on GitHub APIs to extract commits, pull requests, and associated discussions. Although this ensures reproducibility, platform limitations such as incomplete metadata and the absence of external communication channels may affect the full reconstruction of development context.

Since this study uses thematic analysis, there is an inherent risk of interpretative bias, as different researchers may code the same data in different ways. To mitigate this risk, we adopted independent coding, iterative refinement of the codebook, and consensus-based discussion, although some degree of subjectivity remains unavoidable due to the qualitative nature of the analysis.

Additionally, the survey introduces self-selection and response bias, as participation was voluntary and only 20 out of 117 contacted developers responded (17.09\% response rate). This may limit the representativeness of the sample, as responses may be skewed toward developers more engaged with or interested in Type Hint modernization. Recall bias may also affect how developers report past decisions, potentially influencing response accuracy.

Finally, the results reflect a specific temporal window of the Python ecosystem. Given the rapid evolution of development tools, especially with the increasing adoption of LLM-based systems, the findings should be interpreted as representative of the observed period rather than a stable characterization of current or future practices.
